\documentclass[11pt,a4paper]{article}

\usepackage[T1]{fontenc}
\usepackage[utf8]{inputenc}
\usepackage{amsmath,amssymb}
\usepackage{siunitx}
\usepackage{booktabs}
\usepackage{array}
\usepackage{enumitem}
\usepackage{xurl}
\usepackage[hidelinks]{hyperref}
\usepackage[margin=2.1cm]{geometry}
\usepackage[most]{tcolorbox}

\sisetup{detect-all=true,per-mode=symbol}
\setlist[itemize]{leftmargin=*,nosep}
\setlist[enumerate]{leftmargin=*,nosep}
\renewcommand{\arraystretch}{1.18}

\newcommand{\dd}{\Delta}
\newcommand{\kb}{k_{\mathrm{B}}}
\newcommand{\VT}{V_{\mathrm{T}}}
\newcommand{\nint}{n_i}

\title{Carrier Generation and Recombination\\ELEC4603 Wiki Reference}
\author{Prepared for Group 13 wiki revision}
\date{\today}

\begin{document}
\maketitle

\begin{abstract}
This document consolidates the wiki material on carrier generation, recombination, carrier lifetime, diffusion length, device applications, and worked sample problems. It is intended as a paste-ready technical reference for OneNote and as a shareable companion to the interactive demonstrations.
\end{abstract}

\tableofcontents
\newpage

\section{Key Definitions}

Carrier generation creates mobile electron-hole pairs by promoting electrons from the valence band to the conduction band. Carrier recombination removes electron-hole pairs when conduction-band electrons lose energy and fill valence-band hole states. At thermal equilibrium, microscopic generation and recombination still occur, but their rates balance and the carrier concentrations remain constant \cite{entner2007,sze2007}.

\begin{equation}
  n_0p_0=\nint^2 .
\end{equation}

When a semiconductor is disturbed by light, voltage bias, heating, or injection,
\begin{equation}
  n=n_0+\dd n,\qquad p=p_0+\dd p .
\end{equation}
For electron-hole pair generation, the excess carriers are generated together, so \(\dd n=\dd p\). Throughout this document the sign convention is
\begin{equation}
  U=R-G,
\end{equation}
so \(U>0\) means net recombination and \(U<0\) means net generation.

\section{Generation Mechanisms}

Generation occurs when energy creates an electron-hole pair. The dominant mechanism depends on temperature, bandgap, illumination, defect density, and electric field strength \cite{entner2007,sze2007}.

\begin{table}[h]
\centering
\caption{Important carrier generation mechanisms.}
\begin{tabular}{p{0.22\linewidth}p{0.24\linewidth}p{0.24\linewidth}p{0.22\linewidth}}
\toprule
Mechanism & Energy source & Important when & Main device effect\\
\midrule
Thermal & Heat / lattice vibrations & High temperature, small bandgap & Leakage and dark current\\
Optical & Photons & \(h\nu\ge E_g\) & Photocurrent\\
Trap-assisted & Defect levels & Depletion regions and poor material & Reverse leakage\\
Impact ionisation & High electric field & Large reverse bias & Avalanche multiplication\\
\bottomrule
\end{tabular}
\end{table}

\subsection{Thermal Generation}

Thermal generation occurs because the lattice has thermal energy at finite temperature:
\begin{equation}
  \mathrm{thermal\ energy}\rightarrow e^-+h^+ .
\end{equation}
It is strongly connected to intrinsic carrier concentration,
\begin{equation}
  \nint \propto T^{3/2}\exp\left(-\frac{E_g}{2\kb T}\right).
\end{equation}
Higher temperature increases thermal generation, while a larger bandgap reduces it. Thermal generation is often unwanted because it contributes to reverse leakage in diodes and dark current in photodiodes and image sensors \cite{pveducation_i0,mit6720}.

\subsection{Optical Generation}

Optical generation occurs when an absorbed photon has enough energy to excite an electron across the bandgap:
\begin{equation}
  h\nu \ge E_g,\qquad \frac{hc}{\lambda}\ge E_g .
\end{equation}
If the photon energy is below the bandgap, ideal band-to-band absorption cannot create a mobile electron-hole pair. If photon energy is above the bandgap, absorption can generate a mobile electron and hole \cite{entner2007,pveducation_diffusion}.

For a photon flux \(\Phi_0\), a simple depth-dependent generation model is
\begin{equation}
  \Phi(x)=\Phi_0e^{-\alpha x},\qquad
  G_{\mathrm{opt}}(x)=\eta\alpha\Phi_0e^{-\alpha x}.
\end{equation}
If the incident optical intensity is \(I_0\), then \(\Phi_0\approx I_0/(h\nu)\), so
\begin{equation}
  G_{\mathrm{opt}}(x)=\eta\alpha\frac{I_0}{h\nu}e^{-\alpha x}.
\end{equation}
Large absorption coefficient \(\alpha\) means near-surface generation; smaller \(\alpha\) means deeper penetration.

\subsection{Trap-Assisted Generation}

Defects, impurities, and interface states create energy levels inside the bandgap. These levels can act as recombination-generation centres. In a depletion region, mobile carrier concentrations are low, so trap emission can generate carriers that are immediately swept apart by the electric field. A useful approximate depletion-region rate is
\begin{equation}
  G_{\mathrm{SRH}}\simeq \frac{\nint}{\tau_g},
\end{equation}
where \(\tau_g\) is the generation lifetime \cite{shockley1952,hall1952,sze2007}.

\subsection{Impact Ionisation}

Impact ionisation occurs when a high electric field accelerates a carrier so strongly that a collision creates an additional electron-hole pair:
\begin{equation}
  \mathrm{fast\ carrier}+\mathrm{valence\ electron}
  \rightarrow \mathrm{extra\ electron-hole\ pair}.
\end{equation}
The new carriers can also be accelerated, producing avalanche multiplication. This is harmful in ordinary diodes if uncontrolled, but useful in avalanche photodiodes \cite{entner2007,sze2007}.

\section{Recombination Mechanisms}

Recombination mechanisms differ mainly by where the released energy goes and whether defects or extra carriers are involved \cite{entner2007,sze2007}.

\begin{table}[h]
\centering
\caption{Important carrier recombination mechanisms.}
\begin{tabular}{p{0.25\linewidth}p{0.25\linewidth}p{0.25\linewidth}p{0.17\linewidth}}
\toprule
Mechanism & Released energy goes to & Dominant when & Effect\\
\midrule
Radiative & Photon & Direct-bandgap materials & LEDs, lasers\\
Shockley-Read-Hall & Lattice through traps & Defects/traps present & Lifetime and leakage loss\\
Auger & Another carrier & High doping or injection & Efficiency loss\\
Surface & Interface traps & Poor passivation & Optical and interface losses\\
\bottomrule
\end{tabular}
\end{table}

\subsection{Radiative Recombination}

Radiative recombination is direct band-to-band recombination:
\begin{equation}
  e^-+h^+\rightarrow h\nu .
\end{equation}
The common net rate is
\begin{equation}
  U_{\mathrm{rad}}=B(np-\nint^2),
\end{equation}
where \(B\) is the radiative coefficient. Direct-bandgap materials such as GaAs support efficient radiative recombination because energy and momentum conservation are easier than in indirect silicon \cite{entner2007,nptel_band}.

\subsection{Shockley-Read-Hall Recombination}

Shockley-Read-Hall (SRH) recombination occurs through trap states inside the bandgap. The standard net recombination rate is \cite{shockley1952,hall1952,sze2007}
\begin{equation}
  U_{\mathrm{SRH}}=
  \frac{np-\nint^2}
  {\tau_{p0}(n+n_1)+\tau_{n0}(p+p_1)} ,
\end{equation}
where
\begin{equation}
  n_1=\nint\exp\left(\frac{E_t-E_i}{\kb T}\right),\qquad
  p_1=\nint\exp\left(\frac{E_i-E_t}{\kb T}\right).
\end{equation}
Midgap traps are usually strong recombination centres because they are well positioned to capture both electrons and holes.

\subsection{Auger Recombination}

Auger recombination transfers the recombination energy to a third carrier rather than to a photon:
\begin{equation}
  U_{\mathrm{Auger}}=(C_n n+C_p p)(np-\nint^2),
\end{equation}
where \(C_n\) and \(C_p\) are Auger coefficients. It becomes important under high carrier concentration, high-level injection, or heavy doping \cite{entner2007,sze2007}.

\subsection{Surface Recombination}

Surface recombination occurs at semiconductor surfaces and interfaces. Dangling bonds and interface traps provide recombination paths similar to SRH recombination. A common surface recombination flux model is
\begin{equation}
  R_s=S\dd n_s,
\end{equation}
where \(S\) is the surface recombination velocity and \(\dd n_s\) is excess carrier concentration at the surface. A diffusion boundary condition form is
\begin{equation}
  D_n\left.\frac{d\dd n}{dx}\right|_{\mathrm{surface}}=S_n\dd n_s .
\end{equation}

\subsection{Effective Lifetime}

When independent recombination mechanisms act at the same time, their rates add:
\begin{equation}
  \frac{1}{\tau_{\mathrm{eff}}}
  =
  \frac{1}{\tau_{\mathrm{rad}}}
  +\frac{1}{\tau_{\mathrm{SRH}}}
  +\frac{1}{\tau_{\mathrm{Auger}}}
  +\frac{1}{\tau_{\mathrm{surface}}}.
\end{equation}
The shortest lifetime usually dominates because it contributes the largest rate.

\section{Carrier Lifetime and Diffusion Length}

Carrier lifetime \(\tau\) describes how long excess carriers survive before recombining. Diffusion length \(L\) describes how far they travel before recombining. These quantities determine whether generated carriers are collected or lost \cite{pveducation_diffusion,mitpv_materials}.

\subsection{Excess Carrier Decay}

Under low-level injection, the net recombination rate for minority electrons in p-type material is
\begin{equation}
  U_n\simeq \frac{\dd n}{\tau_n}.
\end{equation}
For minority holes in n-type material,
\begin{equation}
  U_p\simeq \frac{\dd p}{\tau_p}.
\end{equation}
When the excitation is removed, excess carriers decay exponentially:
\begin{equation}
  \dd n(t)=\dd n(0)e^{-t/\tau_n}.
\end{equation}
After one lifetime,
\begin{equation}
  \dd n(\tau_n)=\dd n(0)e^{-1}\simeq0.368\,\dd n(0).
\end{equation}
Under steady uniform generation,
\begin{equation}
  \dd n_{\mathrm{ss}}=G_L\tau_n.
\end{equation}

\subsection{Diffusion Length}

The diffusion lengths are
\begin{equation}
  L_n=\sqrt{D_n\tau_n},\qquad L_p=\sqrt{D_p\tau_p}.
\end{equation}
The Einstein relation connects diffusivity and mobility:
\begin{equation}
  D_n=\mu_n\frac{\kb T}{q}=\mu_n\VT,\qquad
  D_p=\mu_p\frac{\kb T}{q}=\mu_p\VT .
\end{equation}
A larger diffusion coefficient or longer lifetime gives a longer diffusion length.

\begin{table}[h]
\centering
\caption{Why lifetime and diffusion length matter.}
\begin{tabular}{p{0.25\linewidth}p{0.65\linewidth}}
\toprule
Context & Effect\\
\midrule
Solar cells & Longer \(L\) improves carrier collection.\\
Photodiodes & Carriers must be collected before recombining.\\
LEDs & Recombination is desired only when radiative.\\
Fast diodes & Shorter \(\tau\) reduces stored minority charge.\\
Defective material & More traps reduce \(\tau\) and \(L\).\\
\bottomrule
\end{tabular}
\end{table}

\section{Generation and Recombination in Devices}

\begin{table}[h]
\centering
\caption{Useful and unwanted effects in devices.}
\begin{tabular}{p{0.2\linewidth}p{0.35\linewidth}p{0.35\linewidth}}
\toprule
Device & Useful effect & Unwanted effect\\
\midrule
P-N diode & Recombination supports forward current & Generation causes reverse leakage\\
Solar cell & Light generation creates photocurrent & Recombination reduces collected current\\
Photodiode & Optical generation creates signal & Thermal generation causes dark current\\
LED & Radiative recombination emits light & Non-radiative recombination wastes energy\\
Switching diode & Short lifetime improves turn-off & Too much recombination increases loss\\
BJT & Carrier injection enables gain & Base recombination reduces gain\\
\bottomrule
\end{tabular}
\end{table}

\subsection{P-N Junction Diodes}

Forward bias reduces the junction barrier. Majority carriers cross the junction and become minority carriers on the opposite side:
\begin{equation}
  \mathrm{carrier\ injection}\rightarrow
  \mathrm{minority\ carrier\ diffusion}\rightarrow
  \mathrm{recombination}.
\end{equation}
This injected minority-carrier diffusion and recombination supports forward current. In reverse bias, thermal and trap-assisted generation in or near the depletion region is swept out by the electric field, producing leakage \cite{mit6012,tugraz_diode}.

For a long-base diode, an ideal diffusion-limited saturation current is
\begin{equation}
  I_S=qA\nint^2\left(\frac{D_p}{L_pN_D}+\frac{D_n}{L_nN_A}\right).
\end{equation}

\subsection{Solar Cells and Photodiodes}

In solar cells and photodiodes,
\begin{equation}
  \mathrm{light}\rightarrow\mathrm{electron-hole\ pairs}\rightarrow
  \mathrm{separated/collected\ current}.
\end{equation}
Good performance requires high optical generation, low recombination, long lifetime, long diffusion length, and low surface recombination. Thermal and trap-assisted generation cause unwanted dark current.

\subsection{LEDs, Fast Diodes, and Avalanche Devices}

LEDs require strong radiative recombination and low non-radiative recombination. Fast switching diodes benefit from shorter lifetime because stored charge scales approximately as
\begin{equation}
  Q_{\mathrm{stored}}\propto I_F\tau .
\end{equation}
Avalanche devices use high-field impact ionisation. In ordinary diodes this can be destructive; in avalanche photodiodes it provides gain.

\section{Worked Sample Problems}

\subsection{Problem 1: Excess Carrier Decay}

\begin{tcolorbox}[title={Question}]
A silicon sample has \(\dd n(0)=\SI{1.0e15}{cm^{-3}}\) and \(\tau_n=\SI{5}{\micro s}\). After light is switched off, find \(\dd n\) after \SI{5}{\micro s} and the time for \(\dd n\) to fall to 1\% of its initial value.
\end{tcolorbox}

Using \(\dd n(t)=\dd n(0)e^{-t/\tau_n}\),
\begin{align}
  \dd n(\SI{5}{\micro s})
  &=10^{15}e^{-5/5}
  =\SI{3.68e14}{cm^{-3}}.
\end{align}
For 1\% remaining,
\begin{align}
  e^{-t/\tau_n}&=0.01,\\
  t&=\tau_n\ln(100)=5(4.605)=\SI{23.0}{\micro s}.
\end{align}

\subsection{Problem 2: Diffusion Length and Recombination Rate}

\begin{tcolorbox}[title={Question}]
A p-type silicon sample at \SI{300}{K} has \(N_A=\SI{1.0e16}{cm^{-3}}\), \(\nint=\SI{1.0e10}{cm^{-3}}\), \(D_n=\SI{35}{cm^2/s}\), \(\tau_n=\SI{5}{\micro s}\), and \(\dd n=\SI{2.0e14}{cm^{-3}}\). Find \(n_0\), \(L_n\), and \(U_n\).
\end{tcolorbox}

For p-type material, \(p_0\simeq N_A\):
\begin{equation}
  n_0=\frac{\nint^2}{p_0}=\frac{(10^{10})^2}{10^{16}}=\SI{1.0e4}{cm^{-3}}.
\end{equation}
The diffusion length is
\begin{equation}
  L_n=\sqrt{D_n\tau_n}=\sqrt{35(5.0\times10^{-6})}
  =\SI{1.32e-2}{cm}=\SI{132}{\micro m}.
\end{equation}
The recombination rate is
\begin{equation}
  U_n\simeq\frac{\dd n}{\tau_n}
  =\frac{2.0\times10^{14}}{5.0\times10^{-6}}
  =\SI{4.0e19}{cm^{-3}.s^{-1}}.
\end{equation}

\subsection{Problem 3: Reverse Saturation Current}

\begin{tcolorbox}[title={Question}]
A long-base silicon diode at \SI{300}{K} has \(A=\SI{1.0e-4}{cm^2}\), \(\nint=\SI{1.0e10}{cm^{-3}}\), \(N_D=N_A=\SI{1.0e16}{cm^{-3}}\), \(D_p=\SI{12}{cm^2/s}\), \(D_n=\SI{35}{cm^2/s}\), \(L_p=\SI{50}{\micro m}\), and \(L_n=\SI{100}{\micro m}\). Estimate \(I_S\).
\end{tcolorbox}

Use
\begin{equation}
  I_S=qA\nint^2\left(\frac{D_p}{L_pN_D}+\frac{D_n}{L_nN_A}\right).
\end{equation}
Convert lengths: \(L_p=\SI{5.0e-3}{cm}\), \(L_n=\SI{1.0e-2}{cm}\). Then
\begin{align}
  \frac{D_p}{L_pN_D}&=\frac{12}{(5.0\times10^{-3})(10^{16})}=2.4\times10^{-13},\\
  \frac{D_n}{L_nN_A}&=\frac{35}{(1.0\times10^{-2})(10^{16})}=3.5\times10^{-13}.
\end{align}
Since \(qA\nint^2=(1.60\times10^{-19})(10^{-4})(10^{20})=1.60\times10^{-3}\),
\begin{equation}
  I_S=(1.60\times10^{-3})(5.9\times10^{-13})=\SI{9.44e-16}{A}.
\end{equation}

\subsection{Problem 4: Design Comparison}

A short minority-carrier lifetime is useful in a fast switching diode because stored charge must be removed before turn-off:
\begin{equation}
  Q_{\mathrm{stored}}\propto I_F\tau .
\end{equation}
Shorter \(\tau\) means less stored charge and faster turn-off. In a solar cell, however,
\begin{equation}
  L=\sqrt{D\tau}.
\end{equation}
Shorter \(\tau\) reduces diffusion length, so more photogenerated carriers recombine before collection.

\subsection{Problem 5: Surface Recombination Design}

\begin{tcolorbox}[title={Question}]
A thin p-type silicon photodiode has \(D_n=\SI{35}{cm^2/s}\), \(\tau_{\mathrm{bulk}}=\SI{10}{\micro s}\), \(W=\SI{20}{\micro m}\), \(S=\SI{1.0e4}{cm/s}\), and \(G_L=\SI{2.0e20}{cm^{-3}.s^{-1}}\). Estimate \(\tau_{\mathrm{eff}}\), \(L_{\mathrm{eff}}\), and \(\dd n_{\mathrm{ss}}\).
\end{tcolorbox}

For two recombining surfaces,
\begin{equation}
  \frac{1}{\tau_{\mathrm{eff}}}\simeq
  \frac{1}{\tau_{\mathrm{bulk}}}+\frac{2S}{W}.
\end{equation}
With \(W=\SI{2.0e-3}{cm}\) and \(\tau_{\mathrm{bulk}}=\SI{1.0e-5}{s}\),
\begin{equation}
  \frac{1}{\tau_{\mathrm{eff}}}
  =10^5+\frac{2(10^4)}{2.0\times10^{-3}}
  =1.01\times10^7\ \mathrm{s^{-1}},
\end{equation}
so
\begin{equation}
  \tau_{\mathrm{eff}}=\SI{99}{ns}.
\end{equation}
Then
\begin{equation}
  L_{\mathrm{eff}}=\sqrt{35(9.90\times10^{-8})}
  =\SI{18.6}{\micro m},
\end{equation}
and
\begin{equation}
  \dd n_{\mathrm{ss}}=G_L\tau_{\mathrm{eff}}
  =(2.0\times10^{20})(9.90\times10^{-8})
  =\SI{1.98e13}{cm^{-3}}.
\end{equation}
The bulk lifetime was \SI{10}{\micro s}, but surface recombination reduces the effective lifetime to about \SI{99}{ns}, showing why surface passivation is critical.

\section{Interactive Demonstrations}

Two original standalone HTML demonstrations accompany this document.
\begin{itemize}
  \item \textbf{Lifetime and diffusion length simulator}: varies lifetime, surface recombination, generation rate, temperature, and thickness to show excess-carrier decay and spatial survival.
  \item \textbf{P-N diode and photodiode simulator}: varies temperature, doping, lifetime, diffusion constants, area, and illumination to show \(I_S\), photocurrent, \(V_{OC}\), dark I-V, and illuminated I-V.
\end{itemize}

These can be shared through GitHub Pages or opened locally from the repository.

\begin{thebibliography}{10}

\bibitem{entner2007}
R. Entner, ``2.3 Carrier Generation and Recombination,'' TU Wien, 2007. [Online]. Available: \url{https://www.iue.tuwien.ac.at/phd/entner/node11.html}

\bibitem{sze2007}
S. M. Sze and K. K. Ng, \emph{Physics of Semiconductor Devices}, 3rd ed. Hoboken, NJ, USA: Wiley, 2007.

\bibitem{pveducation_i0}
C. Honsberg and S. Bowden, ``Ideality Factor and I0,'' PV Education. [Online]. Available: \url{https://www.pveducation.org/pvcdrom/idealty-factor-and-i0}

\bibitem{mit6720}
J. A. del Alamo, ``Lecture 4: Carrier generation and recombination,'' MIT OpenCourseWare, 6.720J Integrated Microelectronic Devices, Spring 2007. [Online]. Available: \url{https://ocw.mit.edu/courses/6-720j-integrated-microelectronic-devices-spring-2007/6ea95228c96b5cb849eca384b52f4ccf_lecture4.pdf}

\bibitem{pveducation_diffusion}
C. Honsberg and S. Bowden, ``Diffusion Length,'' PV Education. [Online]. Available: \url{https://www.pveducation.org/pvcdrom/pn-junctions/diffusion-length}

\bibitem{shockley1952}
W. Shockley and W. T. Read, ``Statistics of the recombinations of holes and electrons,'' \emph{Physical Review}, vol. 87, no. 5, pp. 835--842, 1952, doi: 10.1103/PhysRev.87.835.

\bibitem{hall1952}
R. N. Hall, ``Electron-hole recombination in germanium,'' \emph{Physical Review}, vol. 87, no. 2, p. 387, 1952, doi: 10.1103/PhysRev.87.387.

\bibitem{nptel_band}
NPTEL, ``Band to Band recombination,'' High Speed Semiconductor Devices archive. [Online]. Available: \url{https://archive.nptel.ac.in/content/storage2/courses/117104071/topic19/19_3.html}

\bibitem{mit6012}
C. G. Fonstad, ``Lecture 6: p-n Junctions: I-V Relationship,'' MIT OpenCourseWare, 6.012 Microelectronic Devices and Circuits, Fall 2009. [Online]. Available: \url{https://ocw.mit.edu/courses/6-012-microelectronic-devices-and-circuits-fall-2009/a342ce0afb84bfa8e5b7dc819ff3e9bc_MIT6_012F09_lec06.pdf}

\bibitem{tugraz_diode}
P. Hadley, ``PN diode IV characteristics,'' TU Graz. [Online]. Available: \url{https://lampz.tugraz.at/~hadley/psd/L6/pnIV.php}

\bibitem{mitpv_materials}
T. Buonassisi, ``Lecture 8: Materials Parameters,'' MIT OpenCourseWare, 2.627 Fundamentals of Photovoltaics, Fall 2013. [Online]. Available: \url{https://ocw.mit.edu/courses/2-627-fundamentals-of-photovoltaics-fall-2013/9b4a45fd2ca234fa5f3184f24a991020_MIT2_627F13_lec08.pdf}

\end{thebibliography}

\end{document}
